% Human source transcription, not native author TeX.
% https://ocw.mit.edu/courses/18-712-introduction-to-representation-theory-fall-2010/cabe7aec8816cbba978080c2f1cf5e82_MIT18_712F10_ch5.pdf
% Theorem 5.23；支撑为Lemma 5.16、Definitions 5.20--5.27、Propositions 5.28--5.31、Lemma 5.33及§5.8的全部三个结论。PDF第10--19页。只计Dynkin型方向，不另称证明了非Dynkin型的逆否方向。

\paragraph{记号摘编（助手，据§5.4--§5.5）.}
Work over $\C$. A representation of a quiver assigns a finite-dimensional vector space $V_i$ to each vertex and a linear map to each arrow. Direct sums and isomorphisms are taken vertexwise, compatibly with the arrow maps. In this item the underlying graph $\Gamma$ is one of the simply-laced Dynkin diagrams $A_n,D_n,E_6,E_7,E_8$. On the lattice with basis $\alpha_1,\ldots,\alpha_n$, the form $B=B_\Gamma$ has matrix
\[
B(\alpha_i,\alpha_i)=2,\qquad
B(\alpha_i,\alpha_j)=
\begin{cases}-1,&i,j\text{ adjacent},\\0,&\text{otherwise}\end{cases}\quad(i\ne j).
\]
The positive definiteness of these Dynkin forms is used as in §5.4. A root is a lattice vector $\alpha$ satisfying $B(\alpha,\alpha)=2$. A positive root has all coefficients nonnegative. The representation supported only at vertex $i$, with a one-dimensional space there, is denoted $\C_{(i)}$.

\begin{theorem}[Gabriel, Theorem 5.23; main problem]
Let $Q$ be a quiver of type $A_n,D_n,E_6,E_7,E_8$. Then $Q$ has finitely many indecomposable representations. Namely, the dimension vector of any indecomposable representation is a positive root (with respect to $B_\Gamma$), and for any positive root $\alpha$ there is exactly one indecomposable representation with dimension vector $\alpha$.
\end{theorem}
Here the dimension vector is $d(V)=(\dim V_1,\ldots,\dim V_n)$, and uniqueness is up to isomorphism.

\subsection*{Roots and reflections: supporting material from §5.4}
\begin{lemma}[Lemma 5.16]
Every root is either positive or negative.
\end{lemma}
\begin{proof}
Suppose $\alpha=\sum_i k_i\alpha_i$ has both positive and negative coefficients. Choose vertices $i,j$ with $k_i>0$, $k_j<0$, of minimal possible distance. Let $i'$ be the next vertex after $i$ on the path towards $j$. Remove the edge joining $i$ to $i'$. Since the Dynkin diagram has no cycles, it breaks into connected graphs $\Gamma_1,\Gamma_2$, with $i\in\Gamma_1$, $j\in\Gamma_2$. Put
\[
\beta=\sum_{m\in\Gamma_1}k_m\alpha_m,\qquad
\gamma=\sum_{m\in\Gamma_2}k_m\alpha_m.
\]
Then $\alpha=\beta+\gamma$, and both $\beta$ and $\gamma$ are nonzero. Thus $B(\beta,\beta)\ge2$ and $B(\gamma,\gamma)\ge2$. Furthermore,
\[
B(\beta,\gamma)=-k_i k_{i'}\ge0,
\]
since $k_i>0$ and $k_{i'}\le0$ by the choice of $i,j$. Therefore
\[
B(\alpha,\alpha)=B(\beta,\beta)+2B(\beta,\gamma)+B(\gamma,\gamma)\ge4,
\]
a contradiction.
\end{proof}
\paragraph{Definition 5.20 and Remark 5.21.}
The reflection in a root is $s_\alpha(v)=v-B(v,\alpha)\alpha$. Write $s_i=s_{\alpha_i}$. It fixes vectors orthogonal to $\alpha$ and sends $\alpha$ to $-\alpha$, so it preserves $B$. The simple reflections generate the Weyl group $W$. Since $B$ is positive definite there are finitely many roots. Each $w\in W$ sends each basis root to a root, so there are only finitely many possible $w$, and $W$ is finite.

\subsection*{Reflection functors: §5.6}
\paragraph{Definitions 5.24--5.27.}
A sink is a vertex into which all adjacent arrows point; a source is one from which all adjacent arrows point. Let $\overline Q_i$ be obtained by reversing all arrows at such a vertex $i$.

For a sink, define $F_i^+:\operatorname{Rep}Q\to\operatorname{Rep}\overline Q_i$ by
\[
(F_i^+V)_k=V_k\quad(k\ne i),\qquad
(F_i^+V)_i=\ker\left(\varphi:\bigoplus_{j\to i}V_j\longrightarrow V_i\right).
\]
All maps stay the same except those now pointing out of $i$. These are the compositions of the inclusion of $\ker\varphi$ into $\bigoplus V_j$ with the projections to $V_k$.

For a source, let $\psi:V_i\to\bigoplus_{i\to j}V_j$ be the canonical map. Define $F_i^-:\operatorname{Rep}Q\to\operatorname{Rep}\overline Q_i$ by
\[
(F_i^-V)_k=V_k\quad(k\ne i),\qquad
(F_i^-V)_i=\operatorname{Coker}\psi=\left(\bigoplus_{i\to j}V_j\right)/\im\psi.
\]
All maps stay the same except those now pointing into $i$, which are the inclusions into the direct sum followed by its natural quotient map.

\begin{proposition}[Proposition 5.28]
Let $V$ be an indecomposable representation. At a sink $i$, either $V=\C_{(i)}$, or $\varphi:\bigoplus_{j\to i}V_j\to V_i$ is surjective. At a source $i$, either $V=\C_{(i)}$, or $\psi:V_i\to\bigoplus_{i\to j}V_j$ is injective.
\end{proposition}
\begin{proof}
For the first statement choose a complement $W$ of $\im\varphi$ in $V_i$. Splitting off $W$ gives a direct sum of the representation supported by $W$ only at $i$, and the representation with $\im\varphi$ at $i$ and the original spaces elsewhere. Since $V$ is indecomposable, one summand must be zero. If the first is zero, $\varphi$ is surjective. If the second is zero, $W$ must be one-dimensional, since otherwise it splits as several one-dimensional representations at $i$. The second statement follows similarly by splitting away $\ker\psi$.
\end{proof}
\paragraph{图式转写说明（助手）.} The preceding direct-sum description transcribes the two summand diagrams in the source; it does not add a new lemma.

\begin{proposition}[Proposition 5.29]
If $V$ is surjective at a sink $i$, then $F_i^-F_i^+V\simeq V$. If $V$ is injective at a source $i$, then $F_i^+F_i^-V\simeq V$.
\end{proposition}
\begin{proof}
In this proof $j\to i$ refers to the original quiver. We establish the first statement, restricting to the spaces, as in the source. Only the $i$-th space needs checking. Let $\varphi:\bigoplus_{j\to i}V_j\to V_i$ be surjective and $K=\ker\varphi$. After $F_i^+$, $V_i$ is replaced by $K$, and the map $\psi:K\to\bigoplus_{j\to i}V_j$ is the inclusion. Applying $F_i^-$ replaces $K$ by
\[
K'=\left(\bigoplus_{j\to i}V_j\right)/\im\psi
=\left(\bigoplus_{j\to i}V_j\right)/\ker\varphi
=\im\varphi=V_i.
\]
This is the homomorphism theorem. The other statement is dual.
\end{proof}

\begin{proposition}[Proposition 5.30]
For an indecomposable $V$, the representations $F_i^+V$ and $F_i^-V$, whenever defined, are either indecomposable or zero.
\end{proposition}
\begin{proof}
We prove the case $F_i^+$. By Proposition 5.28, either $\varphi$ is surjective or $V=\C_{(i)}$. In the latter case $F_i^+V=0$. In the former, suppose $F_i^+V=X\oplus Y$ with $X,Y\ne0$. The representation $F_i^+V$ is injective at $i$, since its maps are projections whose direct sum is the tautological embedding. Thus $X,Y$ are also injective at $i$, and Proposition 5.29 gives
\[
F_i^+F_i^-X=X,\qquad F_i^+F_i^-Y=Y.
\]
In particular $F_i^-X,F_i^-Y\ne0$. But
\[
V=F_i^-F_i^+V=F_i^-X\oplus F_i^-Y,
\]
contradicting indecomposability. The case $F_i^-$ is similar.
\end{proof}

\begin{proposition}[Proposition 5.31]
If $V$ is surjective at a sink $i$, then $d(F_i^+V)=s_i d(V)$. If $V$ is injective at a source $i$, then $d(F_i^-V)=s_i d(V)$.
\end{proposition}
\begin{proof}
Again prove the first statement. With $K=\ker\varphi$, surjectivity gives
\[
\dim K=\sum_{j\to i}\dim V_j-\dim V_i.
\]
Consequently the $i$-th entry of $d(F_i^+V)-d(V)$ is $\sum_{j\to i}\dim V_j-2\dim V_i=-B(d(V),\alpha_i)$; every other entry is zero. Thus
\[
d(F_i^+V)=d(V)-B(d(V),\alpha_i)\alpha_i=s_i d(V).
\]
The source case follows similarly.
\end{proof}

\subsection*{Coxeter elements: §5.7}
For a labeling $1,\ldots,n$ of the vertices set $c=s_1s_2\cdots s_n$.
\begin{lemma}[Lemma 5.33]
If $\beta=\sum_i k_i\alpha_i$, with $k_i\ge0$ and not all zero, then $c^N\beta$ has a strictly negative coefficient for some $N\in\N$.
\end{lemma}
\begin{proof}
Since $W$ is finite, $c^M=1$ for some $M$. We claim
\[
1+c+c^2+\cdots+c^{M-1}=0
\]
as operators. This implies the assertion because $\beta$ has a strictly positive coefficient.
It suffices to show that $1$ is not an eigenvalue of $c$: otherwise a nonzero value $w$ of this sum would satisfy $cw=w$. Suppose $cv=v$, so
\[
s_2\cdots s_nv=s_1v.
\]
Since $s_i$ only changes coordinate $i$, we get $s_1v=v$ and $s_2\cdots s_nv=v$. Repeating gives $s_iv=v$ for all $i$, hence $B(v,\alpha_i)=0$ for all $i$. Nondegeneracy of $B$ gives $v=0$, a contradiction.
\end{proof}

\subsection*{Proof of Gabriel's theorem: §5.8}
Label $Q$ so that $i<j$ whenever $j$ is reachable from $i$. This is possible by repeatedly removing a sink and assigning the greatest unused label. Consider
\[
V^{(0)}=V,\quad V^{(1)}=F_n^+V,\quad
V^{(2)}=F_{n-1}^+F_n^+V,\quad\ldots.
\]
These functors are defined: $n$ is a sink of $Q$, $n-1$ is a sink after reflection at $n$, and so on. After $n$ reflections every arrow has been reversed twice, so $V^{(n)}$ is again a representation of $Q$. Continue the same sequence indefinitely.

\begin{theorem}[Theorem 5.34]
For an indecomposable $V$, there exists $m$ such that $d(V^{(m)})=\alpha_p$ for some $p$.
\end{theorem}
\begin{proof}
If $V^{(i)}$ is surjective at the appropriate vertex $k$, Proposition 5.31 gives
\[
d(V^{(i+1)})=d(F_k^+V^{(i)})=s_k d(V^{(i)}).
\]
Thus as long as surjectivity continues, the dimension vectors are $\cdots s_{n-1}s_nd(V)$. This cannot continue indefinitely by Lemma 5.33, since dimension vectors cannot have negative entries. At the first failure of surjectivity, Proposition 5.30 gives indecomposability and Proposition 5.28 gives $d(V^{(i)})=\alpha_p$.
\end{proof}

\begin{corollary}[Corollary 5.35]
The dimension vector of an indecomposable representation of $Q$ is a positive root.
\end{corollary}
\begin{proof}
Theorem 5.34 gives $s_{i_1}\cdots s_{i_m}d(V)=\alpha_p$. Reflections preserve $B$, so
\[
B(d(V),d(V))=B(\alpha_p,\alpha_p)=2.
\]
\end{proof}

\begin{corollary}[Corollary 5.36]
Two indecomposable representations with the same dimension vector are isomorphic.
\end{corollary}
\begin{proof}
Choose $i$ with $d(V^{(i)})=\alpha_p$. The same reflection sequence for $V'$ gives $d(V'^{(i)})=\alpha_p$, so $V'^{(i)}=V^{(i)}=:V^i$ up to isomorphism. Write the sequence as $F_k^+\cdots F_{n-1}^+F_n^+$. All preceding representations are surjective at the respective vertices. Proposition 5.29 therefore gives
\[
\begin{aligned}
F_n^-F_{n-1}^-\cdots F_k^-V^i
&=F_n^-F_{n-1}^-\cdots F_k^-F_k^+\cdots F_{n-1}^+F_n^+V^{(0)}=V,\\
F_n^-F_{n-1}^-\cdots F_k^-V^i
&=F_n^-F_{n-1}^-\cdots F_k^-F_k^+\cdots F_{n-1}^+F_n^+V'^{(0)}=V'.
\end{aligned}
\]
Thus $V\simeq V'$.
\end{proof}
These two corollaries already give finiteness, since there are finitely many roots.

\begin{corollary}[Corollary 5.37]
For every positive root $\alpha$ there is an indecomposable representation $V$ with $d(V)=\alpha$.
\end{corollary}
\begin{proof}
Consider the sequence $s_n\alpha,s_{n-1}s_n\alpha,\ldots$. By Lemma 5.33 a negative root eventually occurs. Let $\beta$ be the root just before the first negative one. Then $\beta$ is positive and $s_i\beta$ is negative for some $i$. Since $s_i$ only changes coordinate $i$, $\beta=\alpha_i$. Hence
\[
(s_q\cdots s_{n-1}s_n)\alpha=\alpha_i.
\]
Let $\C_{(i)}$ be the representation with dimension vector $\alpha_i$ and define
\[
V=F_n^-F_{n-1}^-\cdots F_q^-\C_{(i)}.
\]
This is an indecomposable representation with $d(V)=\alpha$, as required.
\end{proof}
